% =================================================
% Учительская версия рабочего листа
% =================================================

\worksheettitle
  {Длина отрезка}
  {Учительская версия: ответы и методические комментарии}

\begin{teacherbox}[Цель листа]
  Научить различать отрезок и его длину, правильно снимать показания
  линейки, переводить основные единицы длины, строить отрезок заданной длины
  и применять равенство \(AB=AC+CB\).
\end{teacherbox}

\section{Отрезок и его длина}

\segmentlengthfigure

\begin{propertybox}[Ответ]
  Отрезок — геометрическая \textbf{фигура}, а его длина — \textbf{число} с
  единицей измерения. Запись \(AB=6\text{ см }4\text{ мм}\) сообщает
  \textbf{длину} отрезка \(AB\).
\end{propertybox}

\begin{practicebox}[Ответы к заданию 1]
  \begin{enumerate}[label=\textbf{\arabic*.},itemsep=0.35em]
    \item Неверно.
    \item Верно.
    \item Неверно.
    \item Верно.
  \end{enumerate}
\end{practicebox}

\begin{teacherbox}[Диагностика]
  В речи ученик может сказать «отрезок равен пяти сантиметрам». На начальном
  этапе лучше мягко уточнять: «длина отрезка равна пяти сантиметрам».
\end{teacherbox}

\section{Измеряем отрезок}

\correctmeasurementfigure

\begin{propertybox}[Алгоритм]
  Один конец совмещают с делением \(0\), затем читают деление у второго
  конца. В первом рисунке \(AB=4\text{ см }6\text{ мм}\).
\end{propertybox}

\readrulerpracticefigure

\begin{practicebox}[Ответы к заданиям 2 и 3]
  \begin{enumerate}[label=\alph*),itemsep=0.4em]
    \item \(AB=4\text{ см }6\text{ мм}\);
    \item \(CD=7\text{ см }1\text{ мм}-1\text{ см }3\text{ мм}
      =5\text{ см }8\text{ мм}\).
  \end{enumerate}
\end{practicebox}

\begin{teacherbox}[Ключевая ошибка]
  Некоторые ученики называют показание правого конца длиной отрезка даже
  тогда, когда левый конец расположен не у нуля. Спросите: «Какое расстояние
  уже было на линейке до точки \(C\)?»
\end{teacherbox}

\section{Единицы длины}

\unitsladderfigure

\begin{propertybox}[Ответы]
  \[
    1\text{ см}=10\text{ мм},\quad
    1\text{ дм}=10\text{ см},\quad
    1\text{ м}=100\text{ см},\quad
    1\text{ км}=1000\text{ м}.
  \]
\end{propertybox}

\begin{practicebox}[Ответы к заданию 4]
  \begin{enumerate}[label=\alph*),itemsep=0.45em]
    \item \(79\text{ мм}\);
    \item \(508\text{ см}\);
    \item \(9\text{ м }5\text{ см}\);
    \item \(6045\text{ м}\).
  \end{enumerate}
\end{practicebox}

\begin{warningbox}[Исправление ошибки]
  \(6\text{ км}=6000\text{ м}\), поэтому
  \[
    6\text{ км }45\text{ м}=6000\text{ м}+45\text{ м}=6045\text{ м}.
  \]
\end{warningbox}

\begin{teacherbox}[Методический акцент]
  Не сводите перевод единиц только к механическому «дописыванию нулей».
  Ученик должен каждый раз опираться на соотношение единиц и понимать,
  почему числовое значение стало больше или меньше.
\end{teacherbox}

\clearpage
\section{Строим отрезки}

\begin{propertybox}[Правильный порядок]
  \begin{enumerate}
    \item Отметить первый конец.
    \item Совместить с ним деление \(0\) линейки.
    \item Отметить второй конец у нужного деления.
    \item Соединить точки.
  \end{enumerate}
\end{propertybox}

\studentconstructionanswersfigure

\begin{teacherbox}[Проверка построения]
  При печати файл должен выводиться в масштабе \(100\%\). Однако итоговую
  точность лучше проверять обычной линейкой ученика, а не только сравнением
  с изображением на экране.
\end{teacherbox}

\section{Равные отрезки}

\equalcontrastingsegmentsfigure

\begin{practicebox}[Ответ к заданию 6]
  Равны отрезки \(AB\) и \(CD\), потому что
  \[
    AB=CD=4\text{ см}.
  \]
\end{practicebox}

\section{Точка внутри отрезка}

\partssegmentfigure

\begin{propertybox}[Ответ]
  Если точка \(C\) лежит между \(A\) и \(B\), то
  \[
    AB=AC+CB.
  \]
\end{propertybox}

\partsproblemfigure

\begin{practicebox}[Ответ к заданию 7]
  \[
    CB=AB-AC=9\text{ см }5\text{ мм}-4\text{ см }8\text{ мм}
      =4\text{ см }7\text{ мм}.
  \]
\end{practicebox}

\begin{checkpointbox}[Ответ к заданию 8]
  \[
    MN=12\text{ см}=11\text{ см }10\text{ мм},
  \]
  \[
    KN=MN-MK=11\text{ см }10\text{ мм}-7\text{ см }6\text{ мм}
      =4\text{ см }4\text{ мм}.
  \]
\end{checkpointbox}

\Needspace{9\baselineskip}
\begin{teacherbox}[Альтернативное вычисление]
  Можно перевести обе длины в миллиметры:
  \[
    120\text{ мм}-76\text{ мм}=44\text{ мм}=4\text{ см }4\text{ мм}.
  \]
  Оба способа полезно сопоставить.
\end{teacherbox}

\section{Типичные ошибки}

\begin{center}
  \renewcommand{\arraystretch}{1.42}
  \begin{tabularx}{\textwidth}{
    >{\raggedright\arraybackslash}p{0.31\textwidth}
    >{\raggedright\arraybackslash}p{0.30\textwidth}
    >{\raggedright\arraybackslash}X}
    \toprule
    \rowcolor{TableHeader}
    Ошибка & Что не понято & Наводящий вопрос \\
    \midrule
    Отсчёт ведётся от края линейки
      & Ноль принят за физический край инструмента
      & Где находится штрих с числом \(0\)? \\
    \addlinespace[0.2em]
    Показание правого конца сразу названо длиной
      & Не учтено начальное показание
      & С какого деления начинается отрезок? \\
    \addlinespace[0.2em]
    \(6\text{ км }45\text{ м}=645\text{ м}\)
      & Не понято, что \(1\text{ км}=1000\text{ м}\)
      & Сколько метров содержится в шести километрах? \\
    \addlinespace[0.2em]
    При нахождении части длины складываются
      & Не различены целое и часть
      & Ищем весь отрезок или недостающую часть? \\
    \bottomrule
  \end{tabularx}
\end{center}
